% 3-2-1-Backup-Regel: 3 Kopien, auf 2 verschiedenen Medien, davon 1 ausser Haus.
% Wird im Skript (Abschnitt Backup) und auf den RAID-Folien verwendet.
\begin{tikzpicture}[
        copy/.style={draw, rounded corners, thick, minimum width=3.2cm, minimum height=1.5cm, align=center, font=\small},
        zone/.style={draw, dashed, rounded corners, inner sep=8pt},
        lbl/.style={font=\footnotesize\bfseries, anchor=south west}
    ]
    % Zuhause: Original und erste Kopie auf einem anderen Medium
    \node[copy, fill=blue!10]   (orig)  at (0,0)   {\faLaptop\\[2pt] Original\\ \footnotesize Laptop (SSD)};
    \node[copy, fill=green!12]  (copy1) at (4,0)   {\faHdd\\[2pt] Kopie 1\\ \footnotesize externe Festplatte / \acrshort{NAS}};
    % Ausser Haus: zweite Kopie
    \node[copy, fill=orange!15] (copy2) at (9.8,0)   {\faCloud\\[2pt] Kopie 2\\ \footnotesize Cloud / Festplatte bei Verwandten};

    \node[zone, fit=(orig)(copy1)]  (home) {};
    \node[zone, fit=(copy2)]        (away) {};
    \node[lbl] at (home.north west) {Zuhause};
    \node[lbl] at (away.north west) {Ausser Haus};

    \draw[-latex, thick] (orig) -- (copy1);
    \draw[-latex, thick] (copy1) -- (copy2);

    % Die drei Regeln unter der Grafik
    \node[font=\small, align=center] at (4.9,-1.9)
    {\textbf{3} Kopien \quad\textperiodcentered\quad auf \textbf{2} verschiedenen Medien \quad\textperiodcentered\quad davon \textbf{1} ausser Haus};
\end{tikzpicture}
